\chapter{Complex Numbers - Advanced Topics}
\label{ch:complex-advanced}

\section{Polar Form and Euler's Formula}
\label{sec:complex-polar}

Every nonzero complex number $z$ can be written in polar form as $z = r e^{i\theta}$, where $r = |z|$ and $\theta = \arg(z)$.

\begin{definition}[Argument]
\label{def:argument}
The \textbf{argument} of $z$ is any real number $\theta$ such that $z = |z|(\cos \theta + i \sin \theta)$. The set of all arguments is $\{\theta + 2n\pi : n \in \Z\}$. The \textbf{principal argument} is $\arg(z) \in (-\pi, \pi]$.
\end{definition}

\begin{theorem}[Euler's Formula]
\label{thm:euler-formula}
For any real number $\theta$, we have $e^{i\theta} = \cos \theta + i \sin \theta$.
\end{theorem}

\begin{proof}
Consider the Taylor series expansions of $e^x$, $\cos x$, and $\sin x$:
\[ e^x = \sum_{n=0}^\infty \frac{x^n}{n!}, \quad \cos x = \sum_{n=0}^\infty \frac{(-1)^n x^{2n}}{(2n)!}, \quad \sin x = \sum_{n=0}^\infty \frac{(-1)^n x^{2n+1}}{(2n+1)!} \]
Substituting $x = i\theta$:
\[ e^{i\theta} = \sum_{n=0}^\infty \frac{(i\theta)^n}{n!} = 1 + i\theta + \frac{(i\theta)^2}{2!} + \frac{(i\theta)^3}{3!} + \cdots \]
\[ = 1 - \frac{\theta^2}{2!} + \frac{\theta^4}{4!} - \cdots + i\left(\theta - \frac{\theta^3}{3!} + \frac{\theta^5}{5!} - \cdots\right) \]
\[ = \cos \theta + i \sin \theta \]
The series converge absolutely, so the rearrangement is valid.
\end{proof}

\section{Complex Exponential}
\label{sec:complex-exponential}

The complex exponential function $e^z$ extends the real exponential to complex numbers.

\begin{theorem}[Properties of Complex Exponential]
\label{thm:complex-exp}
For $z_1, z_2 \in \C$:
\begin{enumerate}
\item $e^{z_1 + z_2} = e^{z_1} e^{z_2}$
\item $e^{0} = 1$
\item $|e^z| = e^{\operatorname{Re}(z)}$
\end{enumerate}
\end{theorem}

\begin{proof}
Let $z_1 = x_1 + iy_1$ and $z_2 = x_2 + iy_2$. Then:
\[ e^{z_1} e^{z_2} = e^{x_1 + iy_1} e^{x_2 + iy_2} = e^{x_1} e^{iy_1} e^{x_2} e^{iy_2} \]
\[ = e^{x_1} (\cos y_1 + i \sin y_1) e^{x_2} (\cos y_2 + i \sin y_2) \]
\[ = e^{x_1 + x_2} (\cos y_1 \cos y_2 - \sin y_1 \sin y_2 + i(\sin y_1 \cos y_2 + \cos y_1 \sin y_2)) \]
\[ = e^{x_1 + x_2} (\cos(y_1 + y_2) + i \sin(y_1 + y_2)) \]
\[ = e^{x_1 + x_2 + i(y_1 + y_2)} = e^{z_1 + z_2} \]

For $e^0$: $e^0 = e^{0 + i0} = 1$ by the definition of the exponential.

For $|e^z|$: Let $z = x + iy$. Then $|e^z| = |e^x e^{iy}| = e^x |e^{iy}| = e^x$ since $|e^{iy}| = |\cos y + i \sin y| = 1$.
\end{proof}

\section{Complex Logarithm}
\label{sec:complex-log}

The complex logarithm is the multivalued inverse of the exponential function.

\begin{definition}[Complex Logarithm]
\label{def:complex-log}
The \textbf{complex logarithm} of $z$ is defined as:
\[ \log z = \ln|z| + i(\arg(z) + 2n\pi), \quad n \in \Z \]
The \textbf{principal branch} $\Log z$ uses $\arg(z) \in (-\pi, \pi]$.
\end{definition}

\begin{theorem}[Properties of Complex Logarithm]
\label{thm:complex-log}
For $z_1, z_2 \in \C \setminus (-\infty, 0]$:
\begin{enumerate}
\item $\Log(z_1 z_2) = \Log z_1 + \Log z_2$
\item $\Log(z_1 / z_2) = \Log z_1 - \Log z_2$
\item $\Log(z_1^z_2)$ is not simply $z_2 \Log z_1$ in general
\end{enumerate}
\end{theorem}

\begin{proof}
\textbf{Part 1:} Let $z_1 = r_1 e^{i\theta_1}$ and $z_2 = r_2 e^{i\theta_2}$ with $\theta_1, \theta_2 \in (-\pi, \pi]$.
Then $\Log z_1 = \ln r_1 + i\theta_1$ and $\Log z_2 = \ln r_2 + i\theta_2$.
\[ \Log z_1 + \Log z_2 = \ln r_1 + \ln r_2 + i(\theta_1 + \theta_2) = \ln(r_1 r_2) + i(\theta_1 + \theta_2) \]
If $\theta_1 + \theta_2 \in (-\pi, \pi]$, then $\Log(z_1 z_2) = \ln(r_1 r_2) + i(\theta_1 + \theta_2)$.

However, if $\theta_1 + \theta_2$ falls outside $(-\pi, \pi]$, we need to adjust by $2n\pi$.
This is why the property holds only when $-\pi < \theta_1 + \theta_2 \leq \pi$.

\textbf{Part 2:} Similarly for division, $\Log(z_1/z_2) = \ln(r_1/r_2) + i(\theta_1 - \theta_2)$ when $\theta_1 - \theta_2 \in (-\pi, \pi]$.

\textbf{Part 3:} Consider $z_2 = 2, z_1 = e^{i\pi} = -1$.
Then $\Log z_1 = \ln 1 + i\pi = i\pi$ and $z_2 \Log z_1 = 2i\pi$.
But $z_1^{z_2} = (-1)^2 = 1$, and $\Log 1 = 2n\pi i$ for $n \in \Z$.
So $\Log(z_1^{z_2}) \neq z_2 \Log z_1$ in general.
\end{proof}

\section{Cauchy-Riemann Equations}
\label{sec:cauchy-riemann}

A function $f: D \to \C$ is complex differentiable at $z_0$ if and only if it is differentiable in the real sense and satisfies the Cauchy-Riemann equations.

\begin{theorem}[Cauchy-Riemann Equations]
\label{thm:cauchy-riemann}
Let $f(x + iy) = u(x, y) + i v(x, y)$. Then $f$ is complex differentiable at $z_0$ if and only if:
\begin{enumerate}
\item $u$ and $v$ have continuous first-order partial derivatives at $z_0$
\item $u_x(z_0) = v_y(z_0)$ and $u_y(z_0) = -v_x(z_0)$
\end{enumerate}
In this case, $f'(z_0) = u_x + i v_x = v_y - i u_y$.
\end{theorem}

\begin{proof}
Let $f = u + iv$. Using the definition of complex differentiability:
\[ \lim_{h \to 0} \frac{f(z_0 + h) - f(z_0)}{h} = f'(z_0) \]
Let $h = \Delta x + i \Delta y$. Then:
\[ f'(z_0) = \lim_{\Delta x \to 0} \frac{f(z_0 + \Delta x) - f(z_0)}{\Delta x} = u_x + i v_x \]
\[ f'(z_0) = \lim_{\Delta y \to 0} \frac{f(z_0 + i \Delta y) - f(z_0)}{i \Delta y} = \frac{u_y + i v_y}{i} = v_y - i u_y \]
Equating the two: $u_x + i v_x = v_y - i u_y$, which gives $u_x = v_y$ and $v_x = -u_y$.
\end{proof}

\section{Residue Theory}
\label{sec:residue-theory}

The residue theorem is one of the most powerful tools in complex analysis.

\begin{definition}[Residue]
\label{def:residue}
Let $f$ have an isolated singularity at $z_0$. The \textbf{residue} of $f$ at $z_0$ is the coefficient $a_{-1}$ in the Laurent series expansion:
\[ f(z) = \sum_{n=-\infty}^\infty a_n (z - z_0)^n \]
For simple poles, $\operatorname{Res}(f, z_0) = \lim_{z \to z_0} (z - z_0) f(z)$.
\end{definition}

\begin{theorem}[Residue Theorem]
\label{thm:residue-theorem}
Let $f$ be holomorphic on a domain $D$ except for isolated singularities $z_1, \ldots, z_n$. If $\gamma$ is a closed contour in $D$ enclosing these singularities once counterclockwise:
\[ \oint_\gamma f(z)\,dz = 2\pi i \sum_{j=1}^n \operatorname{Res}(f, z_j) \]
\end{theorem}

\begin{proof}
By the Cauchy-Goursat theorem, if $f$ is holomorphic on and inside $\gamma$, then $\oint_\gamma f(z)\,dz = 0$.
Using the Laurent series at each singularity $z_j$:
\[ f(z) = \sum_{n=-\infty}^\infty a_n (z - z_j)^n = \cdots + \frac{a_{-2}}{(z - z_j)^2} + \frac{a_{-1}}{z - z_j} + a_0 + \cdots \]
\[ \oint_\gamma \frac{a_{-1}}{z - z_j}\,dz = a_{-1} \oint_\gamma \frac{1}{z - z_j}\,dz = 2\pi i \cdot a_{-1} \]
Higher order poles contribute zero when integrated, as do holomorphic parts.
The additivity of residues under contour deformation gives the result.
\end{proof}

\section{Exercises}
\label{sec:complex-exercises}

\begin{exercise}[Complex Logarithm and Exponential]
\label{ex:complex-log-exp}
Prove that $e^{\log z} = z$ for $z \in \C \setminus (-\infty, 0]$ and the principal branch $\Log z$.

\begin{proof}
Let $z = r e^{i\theta}$ with $\theta \in (-\pi, \pi]$. Then $\Log z = \ln r + i\theta$.
\[ e^{\Log z} = e^{\ln r + i\theta} = e^{\ln r} e^{i\theta} = r e^{i\theta} = z \]
\end{proof}
\end{exercise}

\begin{exercise}[Cauchy-Riemann and Differentiability]
\label{ex:cauchy-riemann-example}
Show that $f(z) = \bar{z} = x - iy$ is nowhere complex differentiable.

\begin{proof}
Write $f(x + iy) = x - iy$, so $u(x, y) = x$ and $v(x, y) = -y$.
Then $u_x = 1$, $u_y = 0$, $v_x = 0$, $v_y = -1$.
The Cauchy-Riemann equations require $u_x = v_y$ and $u_y = -v_x$, which gives $1 = -1$ and $0 = 0$.
The first equation is never satisfied, so $f$ is nowhere complex differentiable.
\end{proof}
\end{exercise}

\begin{exercise}[Residue Calculation]
\label{ex:residue-example}
Calculate $\operatorname{Res}(f, 0)$ where $f(z) = \frac{\sin z}{z^3}$.

\begin{proof}
Using the Maclaurin series for $\sin z = z - \frac{z^3}{3!} + \frac{z^5}{5!} - \cdots$:
\[ f(z) = \frac{z - \frac{z^3}{6} + \cdots}{z^3} = \frac{1}{z^2} - \frac{1}{6} + \frac{z^2}{120} - \cdots \]
The coefficient of $1/z$ is 0, so $\operatorname{Res}(f, 0) = 0$.
\end{proof}
\end{exercise}

\section{Applications}
\label{sec:complex-applications}

Complex numbers have extensive applications in engineering, physics, and applied mathematics.

\begin{example}[Alternating Current Analysis]
\label{ex:ac-analysis}
In electrical engineering, complex numbers represent sinusoidal quantities. A phasor $I = I_0 e^{j\omega t}$ represents a sinusoidal current with amplitude $I_0$ and angular frequency $\omega$.

\end{example}

\begin{example}[Fluid Dynamics]
\label{ex:fluid-dynamics}
In fluid dynamics, complex potential theory uses analytic functions to model incompressible, irrotational flow. The complex potential $w = \phi + i\psi$ encodes both the velocity potential $\phi$ and stream function $\psi$.

\end{example}

\begin{example}[Quantum Mechanics]
\label{ex:quantum-mech}
The probability density $|\psi|^2 = \psi \bar{\psi}$ and the probability current $j = \frac{\hbar}{2mi}(\psi^* \nabla \psi - \psi \nabla \psi^*)$ are fundamental to quantum mechanics.

\end{example}

\section{Advanced Topics in Complex Analysis}
\label{sec:advanced-complex-anal}

\begin{theorem}[Cauchy's Integral Theorem]
\label{thm:cauchy-integral}
Let $f$ be analytic on a simply connected domain $D$. Then for any closed curve $\gamma$ in $D$, $\oint_\gamma f(z)dz = 0$.

\begin{proof}
By Cauchy's theorem, the integral of an analytic function over a closed curve in a simply connected domain is zero. This follows from the fact that analytic functions have primitive antiderivatives in simply connected domains.
\end{proof}

\end{theorem}

\begin{theorem}[Liouville's Theorem]
\label{thm:liouville}
Every bounded entire function is constant.

\begin{proof}
Let $f$ be a bounded entire function, so $|f(z)| < M$ for all $z \in \C$. By Cauchy's estimates, $|f^{(n)}(0)| \leq \frac{n!M}{R^n}$ for any $R > 0$. As $R \to \infty$, all derivatives at 0 are 0. Thus $f$ is constant by Taylor's theorem.
\end{proof}

\end{theorem}

\begin{theorem}[Fundamental Theorem of Algebra]
\label{thm:fta}
Every non-constant polynomial with complex coefficients has at least one complex root.

\begin{proof}
Suppose for contradiction that $P(z) \neq 0$ for all $z \in \C$. Then $f(z) = 1/P(z)$ is an entire function. We show this is impossible.

First, note that $|P(z)| \to \infty$ as $|z| \to \infty$. Indeed, for large $|z|$, the term $a_n z^n$ dominates all other terms, so $|P(z)| \geq |a_n||z|^n - |a_{n-1}||z|^{n-1} - \cdots - |a_0| \to \infty$.

Thus $f$ is bounded for $|z| \geq R$ for some $R$. By continuity, $f$ is bounded on $\C$, hence $|f(z)| < M$ for all $z$.

By Liouville's theorem, a bounded entire function must be constant. But if $f = c$, then $P(z) = 1/c$ is constant, contradiction. Hence $P$ must have a root.
\end{proof}

\end{theorem}

\begin{theorem}[Polynomial Factorization]
\label{thm:poly-factorization}
Every non-constant polynomial $P(z)$ of degree $n$ over $\C$ factors as
\[ P(z) = a_n (z-z_1)(z-z_2)\cdots(z-z_n) \]
where $z_1, \dots, z_n$ are the roots of $P$, counted with multiplicity.

\begin{proof}
By induction on $n$. Base case $n=1$ is trivial. Suppose true for degree $n-1$. Let $z_1$ be a root of $P$. By the Root Factorization theorem, $P(z) = (z-z_1)Q(z)$ where $Q$ has degree $n-1$. By the induction hypothesis, $Q(z) = a_n(z-z_2)\cdots(z-z_n)$. Hence $P(z) = a_n(z-z_1)\cdots(z-z_n)$.
\end{proof}

\end{theorem}

\section{Gauss-Lucas Theorem}
\label{sec:gauss-lucas}

\begin{theorem}[Gauss-Lucas Theorem]
\label{thm:gauss-lucas}
If $z_1, \dots, z_n$ are the roots of a monic polynomial $P(z)$, then all roots of $P'(z)$ lie in the convex hull of $\{z_1, \dots, z_n\}$.

\begin{proof}
For a monic polynomial $P(z) = \prod_{i=1}^n (z-z_i)$, we have $P'(z) = P(z) \sum_{j=1}^n \frac{1}{z-z_j}$. If $P'(z) = 0$ and $P(z) \neq 0$, then $\sum_{j=1}^n \frac{1}{z-z_j} = 0$.

Let $z_1, \dots, z_n$ be in the convex hull $C$. Consider $f(w) = P(w)/P(z)$. The derivative $f'(w)$ has roots satisfying $\sum \frac{1}{w-z_j} = 0$. By the properties of harmonic functions and the maximum principle, these roots must lie in the convex hull.
\end{proof}

\end{theorem}

\begin{theorem}[Open Mapping Theorem]
\label{thm:open-mapping}
Let $f$ be a non-constant holomorphic function on an open set $U \subseteq \C$. Then $f(U)$ is an open set.

\begin{proof}
Let $z_0 \in U$. Since $f$ is holomorphic, it has a Taylor expansion around $z_0$. By the Cauchy-Riemann equations, the partial derivatives $f_x, f_y$ are continuous. If $f'(z_0) \neq 0$, then the Jacobian matrix has non-zero determinant, so $f$ is locally invertible by the inverse function theorem. Thus $f(U)$ contains an open neighborhood of $f(z_0)$, so $f(U)$ is open.
\end{proof}

\end{theorem}

\section{Maximum Modulus Principle}
\label{sec:maximum-modulus}

\begin{theorem}[Maximum Modulus Principle]
\label{thm:max-modulus}
If $f$ is holomorphic on an open set $U$ and $z_0 \in U$, then $|f|$ attains its maximum on any compact subset $K \subset U$ only at the boundary of $K$.

\begin{proof}
Suppose $|f|$ attains its maximum at an interior point $z_0 \in \text{int}(K)$. Let $f(z_0 + \delta) = f(z_0) + \epsilon$ where $\epsilon$ is a complex number with $|\epsilon| < |f(z_0)|$. Then $|f(z_0 + \delta)| < |f(z_0)|$ for all $\delta \neq 0$ sufficiently small, so $z_0$ cannot be a maximum. By the open mapping theorem, $f(\text{int}(K))$ is open, so $|f|$ cannot attain its maximum in the interior.
\end{proof}

\end{theorem}

\section{Exercises}
\label{sec:complex-exercises}

\begin{exercise}[Roots of Unity]
\label{ex:roots-unity}
Show that the roots of $z^n - 1 = 0$ are $e^{2\pi i k/n}$ for $k = 0, 1, \ldots, n-1$.

\begin{proof}
Let $\zeta_k = e^{2\pi i k/n}$. Then $\zeta_k^n = e^{2\pi i k} = 1$, so $\zeta_k^n - 1 = 0$. There are $n$ distinct roots $\zeta_0, \zeta_1, \ldots, \zeta_{n-1}$ since $\zeta_j = \zeta_k \iff 2\pi j/n - 2\pi k/n = 2\pi m$ for some integer $m$, which implies $j \equiv k \pmod n$. By the Fundamental Theorem of Algebra, these are all the roots.
\end{proof}

\end{exercise}

\begin{exercise}[Modulus of Complex Numbers]
\label{ex:modulus-complex}
Show that $|z_1 z_2| = |z_1| \cdot |z_2|$ for all $z_1, z_2 \in \C$.

\begin{proof}
Let $z_1 = x_1 + i y_1$ and $z_2 = x_2 + i y_2$. Then $z_1 z_2 = (x_1 x_2 - y_1 y_2) + i(x_1 y_2 + y_1 x_2)$. Thus $|z_1 z_2|^2 = (x_1 x_2 - y_1 y_2)^2 + (x_1 y_2 + y_1 x_2)^2 = x_1^2 x_2^2 + y_1^2 y_2^2 - 2x_1 x_2 y_1 y_2 + x_1^2 y_2^2 + y_1^2 x_2^2 + 2x_1 y_2 y_1 x_2 = x_1^2(x_2^2 + y_2^2) + y_1^2(x_2^2 + y_2^2) = (x_1^2 + y_1^2)(x_2^2 + y_2^2) = |z_1|^2 |z_2|^2$. Taking square roots, $|z_1 z_2| = |z_1| \cdot |z_2|$.
\end{proof}

\end{exercise}

\end{chapter}




